\chapter{Introduction}\label{ch:introduction}

\pagestyle{fancy}

The medical field is one of the domains where impressive amounts of data can be collected. This data can include both general patient information (age, weight, height, blood type) and clinical records such as treatment plans (prescribed medications, surgical interventions, vaccines), consultation parameters, hospitalizations, allergies and so on. The medical domain holds valuable information which, through proper processing, can help doctors identify \textbf{patterns that repeat across patients}. In this way, patient groups that require similar treatment plans can be determined, and the treatments that worked best for each group can also be observed.

Despite the growing volume of medical data, numerous healthcare facilities globally still lack systems to efficiently store and process patient records. The data either remains written on stacks of paper, or is stored in hospital databases (for hospitals that have a digital database), but remains untouched and unprocessed. The potential of this data is enormous, but wasted in most cases, as Yadav et al. suggest in their work~\cite{Yadav2018MiningEHR}.

In order to implement a system like the one described above, several \textbf{challenges} arise. The \textbf{first challenge} is obtaining and preparing the data that needs to be processed. The \textbf{second challenge} consists in grouping patients using an algorithm that can combine both numerical data (e.g. age, body mass index, quality-adjusted life year) and categorical data (e.g. gender, dominant care plan, condition burden). The \textbf{third challenge} is finding and analyzing clinical event sequences through the use of a sequential pattern mining algorithm. The \textbf{fourth challenge} relates to interpreting the obtained results and querying the clinical database easily, so non-technical users can easily find the information they are searching for.

Implementing the above requires a complete and complex dataset. Considering that real hospital data is difficult to obtain and protected by GDPR regulations, for this thesis the open-source project \textbf{Synthea}~\cite{walonoski2017} was used, through which a CSV dataset of 22,284 patients was generated, each patient having their entire medical history simulated in depth: organizations, providers, encounters, conditions, allergies, observations, immunizations, care plans, procedures, medications. Due to the depth of the relationships between tables, the generated dataset represents a suitable instrument implementing a system for overcoming the obstacles mentioned above.

Since the dataset contains valuable data from both a numerical and a categorical perspective, this thesis proposes the use of \textbf{K-Prototypes} for implementing the clustering process, as this algorithm allows the simultaneous use of both data types. K-Prototypes, introduced by Huang~\cite{huang1998}, extends the well-known K-Means algorithm by adding support for categorical attributes through a combined distance function that uses squared Euclidean distance for numerical features and a simple matching measure for categorical ones.

To find patterns of clinical sequences (e.g., particular diagnosis - medicine - surgery - care plan), the thesis adopts the \textbf{PrefixSpan} approach~\cite{pei2001}, which is one of the methods for mining sequential patterns and makes it easy to search for frequent sequences within a sequence of events. Using the PrefixSpan technique on the sequences of patient’s medical events extracted from Synthea data, the platform may identify common clinical sequences.

In addition to the clustering and sequential analysis, the system utilizes a \textbf{Large Language Model (LLM)} that performs two functions. First, an LLM generates automatic, easily readable summaries for each patient cluster, converting the numeric statistical outputs into something that a physician would be able to comprehend without specialized training. The second function of an LLM is the possibility of querying the database in natural language via a text-to-SQL agent ("Show patients with repeated asthma hospitalizations," for instance), where the results will also appear in natural language.

This thesis introduces \textbf{HospitAI}, which is a hospital analytics platform with functionalities such as patient clustering using the K-Prototypes algorithm, mining of sequential patterns via PrefixSpan, and natural language interaction supported by an LLM. The platform was developed using synthetic EHRs generated via Synthea. It offers three role-based interfaces: one for patients to consult their own medical record, one for doctors to explore patient profiles and enter consultation data, and another for managers to analyze the results of clustering and sequential pattern mining and to view hospital statistics.

The remainder of this document is organized as it follows. Chapter~\ref{ch:objectives} defines the project objectives. Chapter~\ref{ch:research} reviews the related literature and positions HospitAI within the current state of the art. Chapter~\ref{ch:analysis} presents the theoretical foundation, including the algorithms used and the system architecture. Chapter~\ref{ch:design} describes the detailed design and implementation. Chapter~\ref{ch:testing} presents the testing and validation results. Chapter~\ref{ch:manual} provides the user manual. Chapter~\ref{ch:conclusions} concludes the thesis and discusses future work.